% ============================================================
%  Setmana 2 — Unitat 7 · Límits i continuïtat
%  Exercicis: 66, 68, 70, 76, 88, 90
% ============================================================

\setmana{2}{Unitat 7 \textperiodcentered{} Límits i continuïtat}{27 de setembre de 2026}%
        {66, 68, 70, 76, 88 i 90}

\begin{errada}{Exercici 76 f)}
  \ondiu{$=\dlim_{x\to 3}\dfrac{(3-x)(3+x)}{(x-3)(x+2)}
    =\dlim_{x\to 3}\dfrac{3-x}{x+2}=0$}
  \hadedir{$=\dlim_{x\to 3}\dfrac{(3-x)(3+x)}{(x-3)(x+2)}
    =\dlim_{x\to 3}\dfrac{-(3+x)}{x+2}=-\dfrac{6}{5}$}
  \motiu{Els factors que se simplifiquen són $3-x$ i $x-3$, que són oposats:
    $\dfrac{3-x}{x-3}=-1$. Queda $-\dfrac{3+x}{x+2}$, que en $x=3$ val
    $-\dfrac{6}{5}$.}
\end{errada}

%\senseerrades{66, 68, 70, 88 i 90}

% ---------- Observacions (només surten amb \obstrue) ----------
\ifobs
\bigskip
\noindent
Una observació menor:\par\nopagebreak

\begin{observacio}{Exercici 68}
  A la solució, els apartats c) i d) surten en ordre canviat: primer el d) i
  després el c). Els valors de cadascun són correctes.
\end{observacio}
\fi
